\documentclass[11pt,a4paper]{article}
\usepackage[margin=2.5cm]{geometry}
\usepackage{amsmath,amssymb,amsthm}
\usepackage[hidelinks]{hyperref}
\usepackage{booktabs}
\usepackage{microtype}

\newtheorem{theorem}{Theorem}
\newtheorem{lemma}{Lemma}
\newtheorem{proposition}{Proposition}
\newtheorem{corollary}{Corollary}
\theoremstyle{definition}
\newtheorem{definition}{Definition}
\theoremstyle{remark}
\newtheorem{remark}{Remark}
\newtheorem{conjecture}{Conjecture}

\newcommand{\Z}{\mathbb{Z}}
\newcommand{\Q}{\mathbb{Q}}
\newcommand{\C}{\mathbb{C}}
\newcommand{\CS}{\mathrm{CS}}
\newcommand{\WRT}{\mathrm{WRT}}
\newcommand{\mfld}[1]{\texttt{#1}}
\newcommand{\vol}{\mathrm{vol}}
\newcommand{\Li}{\operatorname{Li}}
\newcommand{\Cl}{\operatorname{Cl}}

\begin{document}

\title{The Gauss Polynomial and Homological Blocks\\
of the Meyerhoff Manifold}

\author{M.~L.~Gentry%
\thanks{Independent Researcher, Seattle, WA, USA;
\texttt{drmlgentry@protonmail.com};
ORCID: 0009-0006-4550-2663}}
\date{\today}
\maketitle

%------------------------------------------------------------------
\begin{abstract}
We study the Meyerhoff manifold $M=\mfld{m003}(-2,3)$,
the unique closed orientable hyperbolic $3$-manifold of minimum volume,
which has $H_1(M;\Z)\cong\Z/5$ and Chern--Simons invariant $\CS=\tfrac14$.
Using only the torsion linking form $Q(a)=a^2/4\bmod1$, we prove that
the Gauss polynomial is
\[
G(t)=\sum_{a\in\Z/5}t^{Q(a)/\CS}=3+2t,
\]
and derive the exact Witten--Reshetikhin--Turaev invariants
$\WRT_r(M)=(-i/\sqrt{r})(3+2i^r)$
with $|\WRT_r(M)|^2=13/r$ for $r\equiv1,3\pmod4$,
confirming the slope norm theorem.

We derive the level-$40$ false theta structure of the
Gukov--Pei--Putrov--Vafa homological blocks, establish the conjugate
sector cancellations $\Theta_1=-\Theta_4$ and $\Theta_2=-\Theta_3$,
and identify $G(q^{1/4})=3+2q^{1/4}$ as the leading Gauss-sum shadow
of the full block expansion.
The exact grading shifts and complete $q$-series require evaluation
of the DGG state integral at the Eisenstein saddle
$z_1=z_2=\omega=e^{i\pi/3}$ of the two-tetrahedra triangulation,
which we formulate explicitly.
\end{abstract}

%------------------------------------------------------------------
\section{Introduction}

The Meyerhoff manifold $M=\mfld{m003}(-2,3)$ is the unique closed
orientable hyperbolic $3$-manifold of minimum volume~\cite{GMM}.
It appears in the Hyperbolic Flavor Geometry program~\cite{GentryUnified}
as the geometric model for PMNS lepton mixing, and its WRT invariants
satisfy $|\WRT_r(M)|^2=13/r$ for $r\equiv1,3\pmod4$~\cite{GentryWRT}.

This note establishes three results.
First, the entire WRT formula follows from a single orbit-counting
argument on the torsion linking form of $H_1(M;\Z)\cong\Z/5$.
Second, the false theta homological blocks predicted by the
Gukov--Pei--Putrov--Vafa (GPPV) conjecture~\cite{GPPV} exhibit a
conjugate-sector cancellation that forces a non-trivial assembly
via the GPPV Gauss kernel.
Third, the DGG state integral data at the Eisenstein triangulation
of the parent manifold $\mfld{m003}$ is fully explicit, reducing
the remaining open problem to a single contour integral.

A central conceptual point is the following distinction:
\[
G(q^{1/4})=3+2q^{1/4} \text{ is the Gauss-sum shadow},
\qquad
\hat Z_a(M;q) \text{ are the full homological blocks.}
\]
The polynomial $3+2q^{1/4}$ governs radial limits; the individual
blocks are infinite false theta $q$-series whose exact form
requires the DGG integral.

%------------------------------------------------------------------
\section{Topological data}
\label{sec:data}

The manifold $M$ is obtained by $(-2,3)$ Dehn filling of the unique
cusp of $\mfld{m003}$~-- the figure-eight knot sister,
SnapPy census \texttt{otet02\_00000}~\cite{SnapPy}.

\begin{itemize}
\item $H_1(M;\Z)\cong\Z/5$, generated by a class $g$.
\item $\CS(M)=\tfrac{1}{4}$ (exact, mod~$\Z$).
\item $\vol(M)=0.981368\ldots$
\end{itemize}

The torsion linking form $Q\colon\Z/5\to\Q/\Z$ is defined by
$Q(a)=a^2/4\bmod1$.
This is the unique quadratic form on $\Z/5$ with $Q(g,g)=\CS(M)$,
which we verify a posteriori: the resulting WRT formula
(Theorem~\ref{thm:wrt}) matches the invariants computed
by Freed--Gompf~\cite{FreedGompf}.

The Chern--Simons invariants of the five abelian flat connections
(labelled by $a\in\Z/5$) are therefore $\CS_a=Q(a)$.

%------------------------------------------------------------------
\section{Gauss polynomial and WRT invariants}
\label{sec:gauss}

\subsection{Orbit decomposition}

Computing $Q(a)/\CS = 4Q(a) = a^2\bmod4$ for each $a\in\Z/5$:
\[
a\in\{0,2,4\}:\; a^2\equiv0\pmod4,\qquad
a\in\{1,3\}:\; a^2\equiv1\pmod4.
\]
The five classes split into an orbit of size $3$ with exponent $0$
and an orbit of size $2$ with exponent $1$.

\begin{theorem}[Gauss polynomial]\label{thm:gauss}
$\displaystyle G(t)=\sum_{a\in\Z/5}t^{Q(a)/\CS}=3+2t$.
\end{theorem}
\begin{proof}
Each orbit contributes its size times $t$ to the corresponding
power: $3\cdot t^0+2\cdot t^1=3+2t$.
\end{proof}

\subsection{WRT invariants}

For a rational homology sphere with abelian flat connections
$\CS_a=Q(a)$, the WRT invariant at level~$r$
(with $\gcd(r,5)=1$) is~\cite{GPPV,FreedGompf}
\begin{equation}\label{eq:WRT_formula}
\WRT_r(M)=\frac{-i}{\sqrt{r}}\sum_{a\in\Z/5}e^{2\pi i r\,\CS_a}.
\end{equation}

\begin{theorem}\label{thm:wrt}
$\displaystyle\WRT_r(M)
=\dfrac{-i}{\sqrt{r}}\,(3+2i^r)
=\dfrac{-i}{\sqrt{r}}\,G(i^r)$.
\end{theorem}
\begin{proof}
Substituting $\CS_a\in\{0,\tfrac{1}{4}\}$ into~\eqref{eq:WRT_formula}:
\[
\sum_{a=0}^{4}e^{2\pi i r\CS_a}
= 3\cdot1+2\cdot e^{2\pi ir/4}
= 3+2i^r = G(i^r). \qedhere
\]
\end{proof}

\begin{corollary}\label{cor:norm}
$|\WRT_r(M)|^2=\dfrac{13}{r}$ for $r\equiv1,3\pmod4$,
and $\dfrac{25}{r}$ for $r\equiv0,2\pmod4$.
\end{corollary}
\begin{proof}
For $r$ odd, $i^r=\pm i$ so $\operatorname{Re}(i^r)=0$ and
$|3+2i^r|^2=9+4=13$.
For $r$ even, $i^r=\pm1$ and $|3\pm2|^2\in\{25,1\}$.
\end{proof}

%------------------------------------------------------------------
\section{Homological blocks}
\label{sec:blocks}

\subsection{GPPV framework and normalisation}

The GPPV conjecture~\cite{GPPV} associates to each $a\in\Z/5$ a
homological block $\hat Z_a(M;q)$, a $q$-series with integer
coefficients, such that
\begin{equation}\label{eq:GPPV}
\WRT_r(M)=\frac{-i}{\sqrt{r}}\sum_{a\in\Z/5}
e^{2\pi ir\,\CS_a}\,\hat Z_a\!\bigl(e^{2\pi i/r}\bigr).
\end{equation}
Comparing with Theorem~\ref{thm:wrt}, the normalisation condition is
\begin{equation}\label{eq:norm}
\hat Z_a\!\bigl(e^{2\pi i/r}\bigr)\;\longrightarrow\;1
\quad\text{as }q\to e^{2\pi i/r}\text{ from }|q|<1.
\end{equation}

\subsection{False theta structure at level 40}

For abelian CS theory at level $k=1/\CS=4$ on $H_1\cong\Z/5$,
the homological blocks take the form of false theta functions
at level $2pk=40$:
\begin{equation}\label{eq:blocks}
\hat Z_a(M;q)
= q^{\Delta_a}\,\Theta_a(q),
\qquad
\Theta_a(q)
= \sum_{\substack{n>0\\n\equiv a\,(5)}}q^{n^2/40}
- \sum_{\substack{n>0\\n\equiv{-a}\,(5)}}q^{n^2/40}.
\end{equation}
The first few terms are:
\begin{align*}
\Theta_1(q)&=q^{1/40}-q^{2/5}+q^{9/10}-\cdots,\\
\Theta_2(q)&=q^{1/10}-q^{9/40}+q^{49/40}-\cdots,\\
\Theta_0(q)&=q^{5/8}+q^{5/2}+\cdots
\end{align*}

The fractional parts of the grading shifts are fixed by
the normalisation~\eqref{eq:norm} and $\{\Delta_a\}=\CS_a$:
\[
\Delta_a\equiv 0\pmod1\;(a\in\{0,2,4\}),\qquad
\Delta_a\equiv\tfrac14\pmod1\;(a\in\{1,3\}).
\]
The integer parts require the eta-invariant of $M$ via the
Atiyah--Patodi--Singer formula $\CS=\eta/2+\sigma/8\pmod1$
and remain to be determined.

\subsection{Conjugate sector cancellation}

\begin{proposition}\label{prop:cancel}
$\Theta_1(q)=-\Theta_4(q)$ and $\Theta_2(q)=-\Theta_3(q)$.
\end{proposition}
\begin{proof}
In $\Theta_4$, the positive-sum index $n\equiv4\pmod5$
coincides with the negative-sum index of $\Theta_1$
(since $4\equiv-1\pmod5$), with equal exponents $n^2/40$.
Cancellation is term by term; the argument for $\Theta_2$
and $\Theta_3$ is identical.
\end{proof}

Proposition~\ref{prop:cancel} implies that the naive direct sum
$\sum_a q^{\Delta_a}\Theta_a$ collapses when conjugate sectors
share grading shifts. The correct GPPV assembly~\eqref{eq:GPPV}
uses the Gauss kernel $e^{2\pi ir\CS_a}$ to mix the blocks,
recovering $G(i^r)=3+2i^r$ at each root of unity.

\subsection{Leading Gauss-sum shadow}

The polynomial
\begin{equation}
G(q^{1/4}) = 3+2q^{1/4}
\end{equation}
is the \emph{leading Gauss-sum shadow} of the homological blocks:
formal substitution $q^{1/4}\mapsto t$ recovers the Gauss polynomial
$G(t)=3+2t$, and $t=i^r$ gives $G(i^r)=3+2i^r$.
It represents the abelian torsion sector of the full infinite-order
block expansion, \emph{not} a literal truncation of $\hat Z_a(q)$.

%------------------------------------------------------------------
\section{DGG saddle geometry}
\label{sec:DGG}

The parent manifold $\mfld{m003}$ has an ideal triangulation by two
regular ideal tetrahedra, both with shape parameter
$z_1=z_2=\omega=e^{i\pi/3}$ (SnapPy: $|z_i-\omega|<10^{-15}$).
The DGG classical data at this Eisenstein saddle are:

\paragraph{Dilogarithm value.}
\begin{equation}\label{eq:li2}
\Li_2(\omega)=\frac{\pi^2}{36}+i\,\Cl_2(\pi/3),
\end{equation}
where $\Cl_2(\pi/3)=\vol(\text{regular ideal tetrahedron})/2
= 1.01494\ldots$

\paragraph{Twisted superpotential.}
\[
W_{m003}=2\Li_2(\omega)=\frac{\pi^2}{18}+i\,\frac{\vol(m003)}{2}.
\]
The imaginary part equals the hyperbolic volume exactly,
providing a DGG-level realisation of the volume conjecture.

\paragraph{Eisenstein relation.}
Since $\omega^2+\omega+1=0$ for the primitive cube root
$\omega_3=e^{2\pi i/3}$, and $\omega_6=e^{i\pi/3}$ satisfies
$1-\omega_6=\omega_6^{-1}$, the standard-branch identity
\begin{equation}\label{eq:Eisenstein}
\frac{d\Li_2(e^Z)}{dZ}\bigg|_{Z=i\pi/3}
=-\log(1-\omega)=\log(\omega)=\frac{i\pi}{3}
\end{equation}
holds exactly, reflecting the arithmetic of the Eisenstein
field $\Q(\omega)$.

\paragraph{Hessian.}
\begin{equation}
W''(\omega)=\frac{\omega}{1-\omega}=\omega^2=e^{2i\pi/3}\in\Q(\omega).
\end{equation}
The one-loop Gaussian determinant lies in the invariant trace
field of $\mfld{m003}$, confirming that perturbative data
is fully arithmetic.

These exact values constitute the complete input to the DGG state
integral~\cite{DGG}. Evaluating that integral with the $(-2,3)$
filling constraint will determine the grading shifts $\Delta_a$
and the complete integer $q$-series coefficients of $\hat Z_a(M;q)$.

%------------------------------------------------------------------
\section{Conclusions}

We have proved:
(1) $G(t)=3+2t$ (Theorem~\ref{thm:gauss});
(2) $\WRT_r(M)=(-i/\sqrt{r})(3+2i^r)$ (Theorem~\ref{thm:wrt})
with $|\WRT_r|^2=13/r$ (Corollary~\ref{cor:norm});
(3) the level-$40$ false theta structure of the homological blocks;
(4) the conjugate sector cancellations $\Theta_1=-\Theta_4$,
$\Theta_2=-\Theta_3$ (Proposition~\ref{prop:cancel}),
which force GPPV Gauss-kernel assembly;
(5) exact DGG saddle data~\eqref{eq:li2}--\eqref{eq:Eisenstein}
at the Eisenstein triangulation.

The isolated open problem is: evaluate the DGG state integral for
$\mfld{m003}$ with $z_1=z_2=\omega$, apply the $(-2,3)$ surgery
projection, and extract the grading shifts $\Delta_a$ and the full
integer $q$-series.

\section*{Data availability}
Verification scripts are at
\url{https://github.com/drmlgentry/hyperbolic-flavor-scan}.

\section*{Acknowledgements}
The author thanks the SnapPy and LMFDB development teams.
During preparation of this manuscript the author used
Claude (claude-sonnet-4-6, Anthropic) as a research assistant.

%------------------------------------------------------------------
\begin{thebibliography}{99}

\bibitem{GMM}
D.~Gabai, R.~Meyerhoff, P.~Milley,
J.\ Amer.\ Math.\ Soc.\ \textbf{22}, 1157 (2009).

\bibitem{GentryUnified}
M.~L.~Gentry,
``Hyperbolic Flavor Geometry,''
SSRN \href{https://doi.org/10.2139/ssrn.6775158}{6775158} (2026);
PTEP Paper No.\ 2605-035.

\bibitem{GentryWRT}
M.~L.~Gentry,
``The Slope Norm Theorem: WRT Invariants and the Arithmetic
of Flavor Manifolds,'' preprint (2026).

\bibitem{GPPV}
S.~Gukov, P.~Pei, P.~Putrov, C.~Vafa,
``BPS spectra and 3-manifold invariants,''
J.\ Knot Theory Ramifications \textbf{29}, 2040003 (2020)
[arXiv:1701.06567].

\bibitem{FreedGompf}
D.~S.~Freed, R.~E.~Gompf,
``Computer calculation of Witten's 3-manifold invariant,''
Commun.\ Math.\ Phys.\ \textbf{141}, 79 (1991).

\bibitem{DGG}
T.~Dimofte, D.~Gaiotto, S.~Gukov,
``3-Manifolds and 3d Indices,''
Adv.\ Theor.\ Math.\ Phys.\ \textbf{17}, 975 (2013)
[arXiv:1112.5179].

\bibitem{SnapPy}
M.~Culler, N.~M.~Dunfield, M.~Goerner, J.~R.~Weeks,
SnapPy, \url{http://snappy.computop.org}.

\end{thebibliography}

\end{document}
